\chapter{表达式、作用域和常量}

\section{问题入口：一行表达式为什么读不懂}

初学者经常写出能编译但难读的表达式：

\begin{lstlisting}[language=C++,caption={压缩过度的表达式}]
if (i++ < values.size() && values[i] > limit || force) {
    process(values[i]);
}
\end{lstlisting}

这行代码的问题不只是运算符优先级。它同时修改 \code{i}、访问容器、判断开关，还混合了 \code{&&} 和 \code{||}。读者必须在脑子里模拟执行顺序。更严重的是，\code{i++} 后再访问 \code{values[i]}，访问的是递增后的下标，可能不是作者想要的。

第一条规则：表达式不是越短越好。只要一行表达式同时做“修改状态”和“判断业务”，就值得拆开。

\section{副作用要显式}

副作用指表达式除了计算值，还改变了某些状态。自增、赋值、函数调用都可能有副作用。把副作用拆出来：

\begin{lstlisting}[language=C++,caption={拆开状态推进和判断}]
const std::size_t current = i;
++i;

const bool in_range = current < values.size();
const bool accepted = in_range && values[current] > limit;
if (accepted || force) {
    process(values[current]);
}
\end{lstlisting}

这个版本更长，但每一步都能命名。命名不是废话，它把中间含义固定下来，也让调试器里能直接观察。

\section{作用域限制变量生命}

变量作用域越大，读者越难判断它在哪里被修改。先看反例：

\begin{lstlisting}[language=C++,caption={变量作用域过大}]
int result = 0;
// 很多代码
for (int value : values) {
    result += value;
}
// 又有很多代码使用 result
\end{lstlisting}

如果 \code{result} 只用于求和后立即返回，就让它待在小作用域里：

\begin{lstlisting}[language=C++,caption={把变量限制在使用附近}]
int sum_values(std::span<const int> values)
{
    int result = 0;
    for (int value : values) {
        result += value;
    }
    return result;
}
\end{lstlisting}

函数本身也是作用域工具。把独立计算抽成函数，变量生命周期自然缩短，错误传播范围也缩小。

\section{const 和 constexpr}

\code{const} 表示对象初始化后不能通过这个名字修改：

\begin{lstlisting}[language=C++,caption={const 表达不再修改}]
const int max_retry_count = 3;
\end{lstlisting}

\code{constexpr} 表示值可以在编译期求出：

\begin{lstlisting}[language=C++,caption={编译期常量}]
constexpr int PORT_MIN = 1;
constexpr int PORT_MAX = 65535;
\end{lstlisting}

两者不是同义词。运行时读到的配置值可以是 \code{const}，但不是 \code{constexpr}：

\begin{lstlisting}[language=C++,caption={运行时 const}]
const int port = read_port_from_config();
\end{lstlisting}

它不能被修改，但它的值要运行时才知道。理解这个区别，后面学习数组大小、模板参数和编译期计算会轻松很多。

\section{enum class 防止常量混用}

裸整数常量容易混：

\begin{lstlisting}[language=C++,caption={裸整数状态不清楚}]
constexpr int STATE_OPEN = 1;
constexpr int COLOR_RED = 1;

int state = COLOR_RED; // 编译器无法阻止
\end{lstlisting}

用 \code{enum class}：

\begin{lstlisting}[language=C++,caption={强类型枚举}]
enum class ConnectionState {
    closed,
    connecting,
    open,
};

enum class Color {
    red,
    green,
};
\end{lstlisting}

\code{ConnectionState} 和 \code{Color} 是不同类型，不能随便互相赋值。类型系统替你挡住一类低级错误。

\section{switch 要处理完整状态}

枚举常配合 \code{switch}：

\begin{lstlisting}[language=C++,caption={处理连接状态}]
std::string_view describe(ConnectionState state)
{
    switch (state) {
    case ConnectionState::closed:
        return "closed";
    case ConnectionState::connecting:
        return "connecting";
    case ConnectionState::open:
        return "open";
    }
    return "unknown";
}
\end{lstlisting}

这里保留最后的 \code{return} 是为了满足编译器路径分析。开启警告后，如果以后给枚举增加新状态却忘了处理，编译器通常会提醒。不要随便写 \code{default} 把所有遗漏吞掉；对封闭枚举，显式列出每个状态更利于维护。

\section{完整案例：重试策略}

实现一个重试策略：最多重试三次，只对临时错误重试。

\begin{lstlisting}[language=C++,caption={重试策略}]
enum class ErrorKind {
    none,
    temporary_failure,
    permanent_failure,
};

constexpr int RETRY_MAX_ATTEMPTS = 3;

bool should_retry(ErrorKind error, int attempts)
{
    if (error != ErrorKind::temporary_failure) {
        return false;
    }
    return attempts < RETRY_MAX_ATTEMPTS;
}
\end{lstlisting}

这个小函数看起来简单，但它体现了本章规则：错误类型用枚举表达，魔法数字变成常量，表达式拆成业务判断。测试也很清楚：无错误不重试，永久错误不重试，临时错误在次数内重试，达到上限不重试。

\begin{exercisebox}
把一个复杂 \code{if} 条件拆成三个具名布尔变量。要求变量名表达业务含义，而不是 \code{a}、\code{b}、\code{flag1}。然后给每个条件组合写测试，确认拆分没有改变行为。
\end{exercisebox}

\section{表达式要能被审查}

复杂表达式最大的问题是审查困难。看一个文件过滤条件：

\begin{lstlisting}[language=C++,caption={难以审查的文件过滤条件}]
if (!path.empty() && path.extension() == ".log" && size > 0 &&
    size < max_size || include_empty && path.extension() == ".txt") {
    process(path);
}
\end{lstlisting}

读者必须同时记住优先级：\code{&&} 高于 \code{||}。但业务上真正想表达的可能是：处理非空日志文件，或者在允许空文件时处理文本文件。把它拆成名字：

\begin{lstlisting}[language=C++,caption={用名字表达业务条件}]
const bool is_log_file = path.extension() == ".log";
const bool is_text_file = path.extension() == ".txt";
const bool has_allowed_size = size > 0 && size < max_size;
const bool can_process_log = is_log_file && has_allowed_size;
const bool can_process_empty_text = include_empty && is_text_file;

if (!path.empty() && (can_process_log || can_process_empty_text)) {
    process(path);
}
\end{lstlisting}

拆开后，审查者能逐个检查业务名是否对应实现。括号也让组合关系不依赖读者背优先级。高质量表达式不是最短，而是能被别人可靠审查。

\section{短路求值不能隐藏副作用}

\code{&&} 和 \code{||} 有短路求值：左边已经决定结果时，右边不会执行。这对空指针检查很有用：

\begin{lstlisting}[language=C++,caption={短路保护指针解引用}]
if (user != nullptr && user->active) {
    send_message(*user);
}
\end{lstlisting}

但不要把副作用藏在右侧：

\begin{lstlisting}[language=C++,caption={短路隐藏副作用}]
if (ready || load_cache()) {
    use_cache();
}
\end{lstlisting}

这里 \code{load_cache()} 是否执行取决于 \code{ready}。如果它只是查询还好，如果它会修改全局缓存或写日志，读者很难从条件表达式看出状态变化。更清楚的写法是先显式处理：

\begin{lstlisting}[language=C++,caption={把副作用变成步骤}]
if (!ready) {
    ready = load_cache();
}
if (ready) {
    use_cache();
}
\end{lstlisting}

这就是本章的主线：会改变状态的动作应当像步骤，而不是藏在条件里。

\section{常量命名要表达领域}

常量不是把数字换成大写就结束。比较两个名字：

\begin{lstlisting}[language=C++,caption={弱常量名}]
constexpr int MAX = 3;
\end{lstlisting}

\begin{lstlisting}[language=C++,caption={带领域的常量名}]
constexpr int RETRY_MAX_ATTEMPTS = 3;
\end{lstlisting}

第二个名字告诉读者这是重试次数上限，不是数组长度、端口数量或线程数。大型项目里常量很多，名字必须带上下文。对于本书里的练习，也要从一开始养成这个习惯。

\section{switch 的遗漏要让编译器发现}

前面说不要随便写 \code{default}。再展开一点。假设枚举新增了状态：

\begin{lstlisting}[language=C++,caption={新增状态}]
enum class ConnectionState {
    closed,
    connecting,
    open,
    failed,
};
\end{lstlisting}

如果原来的 \code{switch} 没有 \code{default}，并开启警告，编译器更容易提醒 \code{failed} 没处理。如果你写了 \code{default: return "unknown";}，新增状态可能悄悄走默认路径，测试没覆盖就漏掉。对封闭枚举，显式处理所有枚举值通常更好；对真正来自外部的不可信整数，才需要默认错误路径。

\section{重试策略推导到测试}

重试策略的测试可以直接从业务表推导：

\begin{longtable}{p{0.26\linewidth}p{0.22\linewidth}p{0.22\linewidth}p{0.18\linewidth}}
\toprule
错误 & 已尝试次数 & 上限 & 是否重试 \\
\midrule
\code{none} & \code{0} & \code{3} & 否 \\
\code{permanent_failure} & \code{0} & \code{3} & 否 \\
\code{temporary_failure} & \code{0} & \code{3} & 是 \\
\code{temporary_failure} & \code{2} & \code{3} & 是 \\
\code{temporary_failure} & \code{3} & \code{3} & 否 \\
\bottomrule
\end{longtable}

这个表比“测一下 \code{should_retry}”更具体。它让你看到边界在 \code{attempts < RETRY_MAX_ATTEMPTS}，不是 \code{<=}。很多 off-by-one 错误都能通过这种表提前发现。

\section{本章验收：减少读者脑内执行}

表达式、作用域和常量的共同目标，是减少读者在脑子里模拟程序的负担。变量靠近使用点，状态改变单独成句，业务条件有名字，数字有身份，枚举状态被穷尽处理。做到这些，代码才容易审查、调试和修改。
